\section{Whole-Body Control}
\label{subsec:wbc_dfcip}
The WBC used for TITA turns the reduced-order MPC plan into joint torques by
solving a quadratic program (QP) over the full robot model. Its decision vector
contains generalized accelerations and the three-dimensional forces at the
left and right wheel contacts,
\[
    x = \begin{bmatrix}
        \ddot{q}^{\top} & f_L^{\top} & f_R^{\top}
    \end{bmatrix}^{\top},
    \qquad \ddot{q}\in\mathbb{R}^{n_v},
\]
where $n_v=6+n_j$ counts the six floating-base velocities and the $n_j$
actuated joints ($n_j=8$ for TITA).
The QP computes generalized accelerations and wheel contact forces while
imposing the floating-base equations of motion as equality constraints. The
actuator torques are then calculated from the optimized accelerations and
contact forces using inverse dynamics.\\
The desired WBC references are constructed from the first control input
$u_0$ of the optimized MPC sequence and the current robot state:
CoM position, velocity, and acceleration; left and right wheel trajectories;
and base yaw, angular velocity, and angular acceleration. These references
form the feedforward and PD terms of the commanded task accelerations. Joint
posture is referenced to the nominal configuration, with zero reference joint
velocity and acceleration. The whole-body QP minimizes the following weighted
quadratic cost:
\begin{equation}
\begin{aligned}
H_{\mathrm{WBC}}(x) ={}&
    w_{\mathrm{com}}
    \left\|J_{\mathrm{com}}\ddot{q}
        +\dot{J}_{\mathrm{com}}\dot{q}
        -a_{\mathrm{com}}^{\mathrm{cmd}}\right\|^2
    \\
    &+w_{\mathrm{lw}}
    \left\|J_{\mathrm{lw}}\ddot{q}
        +\dot{J}_{\mathrm{lw}}\dot{q}
        -a_{\mathrm{lw}}^{\mathrm{cmd}}\right\|^2
    +w_{\mathrm{rw}}
    \left\|J_{\mathrm{rw}}\ddot{q}
        +\dot{J}_{\mathrm{rw}}\dot{q}
        -a_{\mathrm{rw}}^{\mathrm{cmd}}\right\|^2
    \\
    &+w_{\mathrm{base}}
    \left\|J_{\mathrm{base}}^{\omega}\ddot{q}
        +\dot{J}_{\mathrm{base}}^{\omega}\dot{q}
        -\alpha_{\mathrm{base}}^{\mathrm{cmd}}\right\|^2
    \\
    &+w_{\mathrm{posture}}
    \left\|S\left(\ddot{q}
    -a_{\mathrm{posture}}^{\mathrm{cmd}}\right)\right\|^2
    +w_{\ddot{q}}\|\ddot{q}\|^2
    +\epsilon_f
        \left(\|f_L\|^2+\|f_R\|^2\right).
\end{aligned}
\label{eq:tita_wbc_cost}
\end{equation}
Here, each commanded task acceleration contains feedforward acceleration and
PD feedback. The matrix $S\in\mathbb{R}^{n_v\times n_v}$ is a diagonal
selection matrix. The complete set of objectives and hard conditions handled by the
TITA whole-body QP is:
\begin{itemize}
    \item tracking CoM translational acceleration and wheels-centre translational acceleration, 
    \item tracking base angular acceleration, 
    \item regulating hip-abduction joint acceleration toward posture $q_0$,
    \item regularizing generalized acceleration $\ddot{q}$;
    \item regularizing contact forces $f_L$ and $f_R$;
    \item enforcing floating-base Newton--Euler dynamics;
    \item enforcing rolling kinematics at the left and right wheels;
    \item constraining predicted joint velocities;
    \item constraining predicted joint positions to their limits;
    \item bounding wheel contact forces by friction pyramids;
    \item enforcing positive normal force at both wheels.
\end{itemize}
In compact form, the constrained QP is
\[
    \begin{aligned}
        \min_x\quad & H_{\mathrm{WBC}}(x),\\
        \text{subject to}\quad & A_{\mathrm{eq}}x=b_{\mathrm{eq}},\\
        & d_{\min}\le Cx\le d_{\max}.
    \end{aligned}
\]
Here, $A_{\mathrm{eq}}$ is the coefficient matrix for the equality constraints
and $b_{\mathrm{eq}}$ is their right-hand-side vector. In the implementation,
these are the variables named \texttt{A\_eq} and \texttt{b\_eq}: they stack
three rolling-constraint equations for each wheel and six floating-base
dynamics equations. Thus, the QP has twelve scalar equality
constraints. The matrix $C$ and bounds $d_{\min},d_{\max}$ collect the
inequality constraints.\\
The joint inequalities constrain the predicted state one WBC sample ahead.
For each actuated joint $j$, the QP uses
\[
\begin{aligned}
    \dot{q}_{j}^{+}
        &= \dot{q}_{j}+\Delta t_{\mathrm{wbc}}\ddot{q}_{j},\\
    q_{j}^{+}
        &= q_{j}+\Delta t_{\mathrm{wbc}}\dot{q}_{j}
           +\frac{\Delta t_{\mathrm{wbc}}^{2}}{2}\ddot{q}_{j},
\end{aligned}
\]
and imposes the hard bounds
\[
    -\dot{q}_{j,\max}\leq \dot{q}_{j}^{+}\leq \dot{q}_{j,\max},
    \qquad
    q_{j,\min}\leq q_{j}^{+}\leq q_{j,\max}.
\]
The position limits are read from the joint ranges in the MuJoCo model. The
implementation sets the velocity bound to $100\,\mathrm{rad/s}$ for every
actuated joint. The inequality block also applies a friction-pyramid bound to
each contact force in its local contact frame and enforces a positive lower
bound on the normal force.\\
The equality constraints encode the physics the QP cannot violate. They
require the wheels to roll without slipping and the accelerations to be
consistent with the floating-base dynamics under the chosen contact forces.
Both conditions are expressed through the wheel kinematics.\\
For each wheel $i\in\{L,R\}$, define the spatial Jacobian of its centre as
\[
    J_i=\begin{bmatrix} J_i^{p} \\ J_i^{\omega} \end{bmatrix}
    \in\mathbb{R}^{6\times n_v}.
\]
Its linear block $J_i^{p}$ is the wheel-centre Jacobian already introduced in the
cost ($J_L^{p}=J_{\mathrm{lw}}$, $J_R^{p}=J_{\mathrm{rw}}$); $J_i^{\omega}$
maps $\dot{q}$ to the angular velocity of the wheel,
$\boldsymbol{\omega}_i=J_i^{\omega}\dot{q}$, expressed in the world frame.
Splitting the columns into floating-base and actuated coordinates gives
\[
    J_i=\begin{bmatrix} J_{iu} & J_{ia} \end{bmatrix},
    \qquad
    J_{iu}\in\mathbb{R}^{6\times 6},\quad
    J_{ia}\in\mathbb{R}^{6\times n_j}.
\]
For the contact point, let $\mathbf{r}_i$ be the
vector from the centre of wheel $i$ to its contact point. It is not read from
the simulator's contacts but built geometrically from a virtual frame
$R_i^{\mathrm{v}}=[\,\mathbf{t}_i\;\mathbf{n}_i\;\mathbf{s}_i\,]$ attached to
the wheel centre: $\mathbf{n}_i$ is the wheel axis, $\mathbf{s}_i$ is the
in-plane direction closest to the world vertical, and
$\mathbf{t}_i=\mathbf{n}_i\times\mathbf{s}_i$ is the rolling direction. In
this frame the contact vector is fixed,
$\mathbf{r}_i=R_i^{\mathrm{v}}[0\;0\;-r_w]^{\top}=-r_w\mathbf{s}_i$, where
$r_w$ is the wheel radius. A flat, horizontal ground is therefore assumed. A contact force $f_i$, expressed in
the world frame, acts at that point; seen from the wheel centre, it becomes a
force plus the moment it generates. The matrix $T_i\in\mathbb{R}^{6\times 3}$
is defined as:
\[
    T_i=\begin{bmatrix} I_3 \\ [\mathbf{r}_i]_{\times} \end{bmatrix},
    \qquad
    T_i f_i
    =
    \begin{bmatrix}
        f_i \\
        \mathbf{r}_i \times f_i
    \end{bmatrix},
\]
where $[\mathbf{r}_i]_{\times}$ is the skew-symmetric matrix such that
$[\mathbf{r}_i]_{\times}f_i=\mathbf{r}_i\times f_i$.\\
The QP imposes the rolling constraint of the wheels at the acceleration
level, with the equation of Kameduła~\cite{kamedula_centauro}:
\[
    \left(J_i^{p}-[\mathbf{r}_i]_{\times}J_i^{\omega}\right)\ddot{q}
    =
    -\left(\dot{J}_i^{p}-[\mathbf{r}_i]_{\times}\dot{J}_i^{\omega}\right)\dot{q}
    -\boldsymbol{\omega}_i\times\left(\boldsymbol{\omega}_i^{\mathrm{v}}\times\mathbf{r}_i\right).
\]
All vectors are expressed in the world frame. Here
$\boldsymbol{\omega}_i^{\mathrm{v}}$ is the angular velocity of the virtual
frame $R_i^{\mathrm{v}}$. The frame follows the wheel when it tilts or steers
but not when it spins, so its angular velocity is the wheel one without the
component along the axis,
\[
    \boldsymbol{\omega}_i^{\mathrm{v}}
    =\left(I_3-\mathbf{n}_i\mathbf{n}_i^{\top}\right)\boldsymbol{\omega}_i .
\]
Regarding the dynamics of the system, the mass matrix $M$ and the bias vector $c$, which gathers Coriolis, centrifugal, and
gravity terms, are split by rows into an unactuated part (subscript $u$, the
first six rows) and an actuated part (subscript $a$). The unactuated rows give
\[
    M_u\ddot{q}+c_u = J_{Lu}^{\top}T_Lf_L+J_{Ru}^{\top}T_Rf_R .
\]
Once $\ddot{q}$ and $f_L,f_R$ have been found, the remaining $n_j$ actuated
rows give the torque command,
\[
    \tau = M_a\ddot{q}+c_a-J_{La}^{\top}T_Lf_L-J_{Ra}^{\top}T_Rf_R.
\]
The Jacobian-transpose products $J_{La}^{\top}T_Lf_L$ and
$J_{Ra}^{\top}T_Rf_R$ map the contact wrenches to generalized forces on the
actuated coordinates. The resulting torques form the actuator-level command
that realizes the motion selected by the whole-body optimization.

\begingroup
\small
\setlength{\tabcolsep}{4pt}
\begingroup
\begin{longtable}{p{0.62\textwidth} p{0.28\textwidth}}
\caption{TITA whole-body controller parameters and QP constraints.}
\label{tab:tita_wbc_parameters}\\
\hline
\textbf{Parameter} & \textbf{Value} \\
\hline
\endfirsthead
\hline
\textbf{Parameter} & \textbf{Value} \\
\hline
\endhead
\hline
\endfoot
\hline
\endlastfoot
\multicolumn{2}{l}{\textbf{Timing and references}}\\
WBC update frequency, $f_{\mathrm{wbc}}$ & 100 Hz\\
WBC sample time, $\Delta t_{\mathrm{wbc}}$ & 0.01 s\\
Reference lookahead, $\Delta t_{\mathrm{lookahead}}$ & 0.0 s\\
Wheel radius, $r_w$ & 0.0925 m\\
\multicolumn{2}{l}{\textbf{PD gains}}\\
CoM/base motion gains & $K_{p,\mathrm{motion}}=50$, $K_{d,\mathrm{motion}}=30$\\
Wheel gains & $K_{p,\mathrm{wheel}}=50$, $K_{d,\mathrm{wheel}}=30$\\
Joint-posture regulation gains & $K_{p,\mathrm{reg}}=100$, $K_{d,\mathrm{reg}}=20$\\
Posture-regulated joints, $\mathcal{J}_{\mathrm{posture}}$ & $\{0,4\}$ (left and right hip-abduction joints)\\
\multicolumn{2}{l}{\textbf{QP objective weights}}\\
CoM acceleration, $w_{\mathrm{com}}$ & 1.0\\
Left-wheel acceleration, $w_{\mathrm{lw}}$ & 1.0\\
Right-wheel acceleration, $w_{\mathrm{rw}}$ & 1.0\\
Base angular acceleration, $w_{\mathrm{base}}$ & 1.0\\
Posture task, $w_{\mathrm{posture}}$ & 0.1\\
Generalized-acceleration regularization, $w_{\ddot{q}}$ & $1\times10^{-12}$\\
Contact-force regularization, $\epsilon_f$ & $1\times10^{-9}$\\
\multicolumn{2}{l}{\textbf{QP constraints and solver}}\\
Friction coefficient, $\mu$ & 0.9\\
Minimum normal contact force, $F_{z,\min}$ & 5.0 N\\
Actuated-joint velocity bound, $\dot{q}_{j,\max}$ & $\pm100$ rad/s\\
Joint-position bounds, $q_{j,\min},q_{j,\max}$ & MuJoCo model joint ranges\\
QP solver & QPAX\\
Solver tolerance, $\epsilon_{\mathrm{QP}}$ & $1\times10^{-3}$\\
\end{longtable}
\endgroup
\endgroup
